\section{总结与延伸}

\subsection{核心结论}

\begin{enumerate}
\item LLM 面临数据边际收益下降，机器人面临真实数据荒漠。
\item 机器人数据应理解为金字塔：真实端侧数据、仿真/合成数据、互联网/人类数据。
\item Data Engine 的核心是持续发现失败、补数据、训练、评测和再部署。
\item 数据价值不按 GB 计算，而按真实性、稀缺性、可验证性和边际收益计算。
\item Data Recipe 是数据产业的隐性壁垒。
\item 机器人产业版图是大模型、世界模型、VLA、本体和仿真数据公司的共生网络。
\item 数据终点不是数据消失，而是从静态数据转向动态环境和自我学习。
\end{enumerate}

\subsection{开放问题}

\begin{itemize}
\item 机器人数据金字塔中，仿真数据能否真正跨越 sim-to-real gap？
\item 谁会掌握机器人数据闭环：本体公司、大模型公司，还是仿真数据公司？
\item Data Recipe 能否成为类似模型架构一样的核心壁垒？
\item 当 AI 自我学习增强，data factory 会变成 environment factory 吗？
\end{itemize}

\subsection{拓展阅读}

\begin{itemize}
\item Tesla FSD data engine：理解端侧数据闭环。
\item Behavior benchmark、VLA、world model 相关论文：理解机器人评测与世界模型关系。
\item 合成数据、仿真、domain randomization、sim-to-real 相关材料：理解机器人数据的核心挑战。
\item EP137 AI for Math：理解可验证环境如何成为高质量数据源。
\end{itemize}

\begin{importantbox}{最后的判断}
数据的未来不是“更多爬虫”，而是“更好的环境”。谁能构造可验证、可规模化、能产生失败反馈的环境，谁就能在下一阶段智能训练中占据主动。
\end{importantbox}

\end{document}
